% !TEX program = latexmk
% !TEX options = -xelatex -synctex=1 -interaction=nonstopmode -file-line-error "%DOC%"
\documentclass[a4paper, 12pt, oneside]{article}

% ================= 基础宏包与环境设定 =================
\usepackage[UTF8]{ctex} % 自动调用宋体、黑体、楷体等
\xeCJKsetup{AutoFakeBold=true} % 允许伪粗体，防止宋体加粗报错
\setCJKmainfont{SimSun}[AutoFakeBold=3,BoldFont=SimHei,ItalicFont=KaiTi]
\setCJKsansfont{SimHei}
\setCJKmonofont{FangSong}
% 确保任何后续创建的 SimSun 字体族都带有粗体映射
\defaultfontfeatures[SimSun]{BoldFont=SimHei,ItalicFont=KaiTi,AutoFakeBold=3}

% 将正文的英文和数字字体设为 Times New Roman
\setmainfont{Times New Roman}

\usepackage[a4paper, top=2.55cm, bottom=2.75cm, left=2.4cm, right=2.4cm, bindingoffset=0.5cm, headheight=15pt, headsep=0.37cm, footskip=0.85cm]{geometry}
\usepackage{setspace} % 行距设置宏包
\setstretch{1.70} % 复刻 Word 1.5 倍行距（LaTeX的1.625 ≈ Word的1.5）

\usepackage{amsmath} % 数学公式宏包
\allowdisplaybreaks[4]
\usepackage{amssymb} % 数学符号宏包
\usepackage{amsthm} % 必须引入的数学定理宏包
\usepackage{newtxmath} % 兼容 XeLaTeX，让公式也是 Times 风格，且不干扰正文新罗马
\usepackage{tikz}             % 绘图基础包
\usepackage{pgfplots}         % 绘制二维/三维坐标图
\pgfplotsset{compat=1.18}     % 设定 pgfplots 版本以避免警告
\usepackage{graphicx} % 插图宏包
\graphicspath{{./}{image/}{../image/}} % 自动搜索当前目录和 image 子目录
\usepackage{booktabs} % 三线表宏包
\usepackage{multirow} % 单元格合并宏包
\usepackage{longtable} % 跨页长表格宏包
\usepackage{tabularx} % 自动调整宽度的表格宏包
\usepackage{array} % 增强型表格列宏包
\usepackage{calc} % 维度算术表达式支持
\providecommand{\real}[1]{#1} % Pandoc 在 longtable 中沿用 tabularx 的 \real，补上通用定义
\usepackage{caption} % 图表标题定制
\usepackage{subcaption} % 子图定制

% ================= 图、表、公式编号与样式自动化 =================
% 规范8：公式编号如“(式1)”，写在右边行末。
\makeatletter
\def\tagform@#1{\maketag@@@{(式#1\unskip\@@italiccorr)}}
\makeatother
\renewcommand{\theequation}{\arabic{equation}} % 内部编号恢复纯数字，防交叉引用套娃

% 规范7：表格采用三线表，上下线用1.5磅粗线，中间线用1磅细线
\setlength{\heavyrulewidth}{1.5pt}
\setlength{\lightrulewidth}{1pt}

% 定义图表标题字体为宋体五号
\DeclareCaptionFont{wuhao}{\zihao{5}\songti}

% 规范6/7：图号/表号后空1个字距写图题/表题
\DeclareCaptionLabelFormat{myfig}{图#2}
\DeclareCaptionLabelFormat{mytab}{表#2}

% 规范6：图题放图下方居中，图表上下各空一行(通过skip调整)
\captionsetup[figure]{labelformat=myfig, labelsep=quad, font=wuhao, position=bottom, skip=1.5ex}

% 规范7：表题放表格上方居中
\captionsetup[table]{labelformat=mytab, labelsep=quad, font=wuhao, position=top, skip=1.5ex}

% 规范6：图中若有分图时，分图号用“(a)”“(b)”置于分图之下
\captionsetup[subfigure]{labelformat=simple, labelsep=space, font=wuhao}
\renewcommand\thesubfigure{(\alph{subfigure})}

% ================= 标题格式与目录自动化 =================
\usepackage{titlesec} % 自定义章节标题
\usepackage{tocloft} % 自定义目录样式

% 规范3/4：标题阿拉伯数字连续编号后空一个字距，无空行
% 一级标题：黑体四号字，左顶格
\titleformat{\section}{\zihao{4}\heiti}{\thesection\quad}{0pt}{}
\titlespacing*{\section}{0pt}{*1.5}{*1.5} % 使用相对单位替代 0pt，恢复正常行距缓冲

% 二级标题：黑体小四号字，左顶格
\titleformat{\subsection}{\zihao{-4}\heiti}{\thesubsection\quad}{0pt}{}
\titlespacing*{\subsection}{0pt}{*1.2}{*1.2}

% 三级标题：楷体小四号字，左顶格
\titleformat{\subsubsection}{\zihao{-4}\kaishu}{\thesubsubsection\quad}{0pt}{}
\titlespacing*{\subsubsection}{0pt}{*1}{*1}

% 四级标题：楷体小四号字，左顶格
\titleformat{\paragraph}{\zihao{-4}\kaishu}{\theparagraph\quad}{0pt}{}
\titlespacing*{\paragraph}{0pt}{*1}{*1}

% 规范目录：目录两字居中，空4个字距，黑体四号；显示三级标题，内容宋体小四
\setcounter{tocdepth}{3}
\renewcommand{\cfttoctitlefont}{\hfill\zihao{4}\heiti}
\renewcommand{\cftaftertoctitle}{\hfill}
\renewcommand{\contentsname}{目\hspace{4em}录}
\renewcommand{\cftsecfont}{\zihao{-4}\songti}
\renewcommand{\cftsecpagefont}{\zihao{-4}\songti}
\renewcommand{\cftsubsecfont}{\zihao{-4}\songti}
\renewcommand{\cftsubsecpagefont}{\zihao{-4}\songti}
\renewcommand{\cftsubsubsecfont}{\zihao{-4}\songti}
\renewcommand{\cftsubsubsecpagefont}{\zihao{-4}\songti}
\renewcommand{\cftsecleader}{\cftdotfill{\cftdotsep}}

% 层次标题序号一律左对齐
\setlength{\cftsecindent}{0pt}
\setlength{\cftsubsecindent}{0pt}
\setlength{\cftsubsubsecindent}{0pt}

% 统一目录行距（增加一定的弹性）
\setlength{\cftbeforesecskip}{0pt plus 5pt}
\setlength{\cftbeforesubsecskip}{0pt plus 3pt}
\setlength{\cftbeforesubsubsecskip}{0pt plus 3pt}

% 调整编号宽度，防止左对齐后长编号（如 1.1.1）与后续汉字粘连
\setlength{\cftsecnumwidth}{1.5em}
\setlength{\cftsubsecnumwidth}{2.2em}
\setlength{\cftsubsubsecnumwidth}{3.0em}

% ================= 页眉页脚与脚注设定 =================
\usepackage{fancyhdr}
% 规范4：全部不加页眉。页码数字两侧不加修饰线，位于页脚居中，五号字。
\fancypagestyle{plain}{
  \fancyhf{}
  \renewcommand{\headrulewidth}{0pt}
  \fancyfoot[C]{\zihao{5}\thepage}
}
\pagestyle{plain}

% 规范5：脚注右上标连续编号，宋体小五号字
\renewcommand{\footnotesize}{\zihao{-5}\songti}

\usepackage[hidelinks]{hyperref} % 提供目录及交叉引用跳转

% ================= 终极交叉引用神器 cleveref (必须放最后) =================
\usepackage{cleveref}

% 汉化并消除英式空格（#2和#3是超链接锚点，#1是数字）
\crefformat{figure}{#2图#1#3}
\crefformat{table}{#2表#1#3}
\crefformat{equation}{#2式#1#3}
\crefformat{theorem}{#2定理#1#3}
\crefformat{lemma}{#2引理#1#3}
\crefformat{remark}{#2注解#1#3}
\crefformat{section}{#2第#1节#3}

% 汉化多重引用时的连词（比如同时引用两个图，三个公式）
\newcommand{\crefpairconjunction}{~和~}
\newcommand{\crefmiddleconjunction}{、}
\newcommand{\creflastconjunction}{~和~}

% ================= 定理、引理与注解环境定制 =================

% 自定义一个适合中文排版的定理样式 "cn_theorem"
\newtheoremstyle{cn_theorem}
  {0pt}           % 段前间距：0pt
  {0pt}           % 段后间距：0pt
  {\normalfont}   % 正文字体：强制使用正常字体（宋体，绝不倾斜）
  {2\ccwd}        % 缩进：首行自动缩进两个中文字符
  {\heiti}        % 表头字体：让“定理”、“引理”显示为黑体
  {}              % 表头后的标点：留空
  {\ccwd}         % 表头与正文之间的空格：空一个中文字符
  {}              % 表头自定义格式：使用默认

% 激活并应用上述样式
\theoremstyle{cn_theorem}

% 【最简单且绝对不会错的方案】：全部独立计数，全挂载在 section 之下！
\newtheorem{theorem}{定理}[section]
\newtheorem{lemma}{引理}[section]   % 不再共享 theorem 的计数器，独立使用自己的
\newtheorem{remark}{注解}[section]  % 注解也独立

% 将默认的证明环境英文 "Proof" 替换为中文黑体的 "证明"
\renewcommand{\proofname}{\heiti 证明}

% ================= 附录自动化命令 =================
\newcommand{\startappendix}{
    \setcounter{section}{0}
    \renewcommand{\thesection}{附录\Alph{section}} 
    \renewcommand{\thefigure}{\Alph{section}\arabic{figure}}
    \renewcommand{\thetable}{\Alph{section}\arabic{table}}
    \renewcommand{\theequation}{式\Alph{section}\arabic{equation}}
    
    \renewcommand{\theHsection}{\Alph{section}}
    \renewcommand{\theHfigure}{\Alph{section}\arabic{figure}}
    \renewcommand{\theHtable}{\Alph{section}\arabic{table}}
    \renewcommand{\theHequation}{\Alph{section}\arabic{equation}}
    
    \setcounter{figure}{0}
    \setcounter{table}{0}
    \setcounter{equation}{0}

    \titleformat{\section}{\zihao{4}\heiti\centering}{\thesection\quad}{0pt}{}
    \addtocontents{toc}{\protect\setlength{\cftsecnumwidth}{4.5em}}
}

% ================= NodePaper generated-template helpers =================
\providecommand{\tightlist}{\setlength{\itemsep}{0pt}\setlength{\parskip}{0pt}}
\usepackage{color}
\usepackage{fancyvrb}
\IfFileExists{fvextra.sty}{
  \usepackage{fvextra}
  \DefineVerbatimEnvironment{Highlighting}{Verbatim}{commandchars=\\\{\},breaklines,breakanywhere,fontsize=\small}
}{
  \DefineVerbatimEnvironment{Highlighting}{Verbatim}{commandchars=\\\{\},fontsize=\small}
}
\newcommand{\VerbBar}{|}
\newcommand{\VERB}{\Verb[commandchars=\\\{\}]}
\usepackage{framed}
\definecolor{shadecolor}{RGB}{248,248,248}
\newenvironment{Shaded}{\begin{snugshade}}{\end{snugshade}}
\usepackage{adjustbox}
\newcommand{\pandocbounded}[1]{\adjustbox{max width=\linewidth,max height=.8\textheight,keepaspectratio}{#1}}
\newcommand{\AlertTok}[1]{\textcolor[rgb]{0.94,0.16,0.16}{#1}}
\newcommand{\AnnotationTok}[1]{\textcolor[rgb]{0.56,0.35,0.01}{\textbf{\textit{#1}}}}
\newcommand{\AttributeTok}[1]{\textcolor[rgb]{0.13,0.29,0.53}{#1}}
\newcommand{\BaseNTok}[1]{\textcolor[rgb]{0.00,0.00,0.81}{#1}}
\newcommand{\BuiltInTok}[1]{#1}
\newcommand{\CharTok}[1]{\textcolor[rgb]{0.31,0.60,0.02}{#1}}
\newcommand{\CommentTok}[1]{\textcolor[rgb]{0.56,0.35,0.01}{\textit{#1}}}
\newcommand{\CommentVarTok}[1]{\textcolor[rgb]{0.56,0.35,0.01}{\textbf{\textit{#1}}}}
\newcommand{\ConstantTok}[1]{\textcolor[rgb]{0.56,0.35,0.01}{#1}}
\newcommand{\ControlFlowTok}[1]{\textcolor[rgb]{0.13,0.29,0.53}{\textbf{#1}}}
\newcommand{\DataTypeTok}[1]{\textcolor[rgb]{0.13,0.29,0.53}{#1}}
\newcommand{\DecValTok}[1]{\textcolor[rgb]{0.00,0.00,0.81}{#1}}
\newcommand{\DocumentationTok}[1]{\textcolor[rgb]{0.56,0.35,0.01}{\textbf{\textit{#1}}}}
\newcommand{\ErrorTok}[1]{\textcolor[rgb]{0.64,0.00,0.00}{\textbf{#1}}
}
\newcommand{\ExtensionTok}[1]{#1}
\newcommand{\FloatTok}[1]{\textcolor[rgb]{0.00,0.00,0.81}{#1}}
\newcommand{\FunctionTok}[1]{\textcolor[rgb]{0.13,0.29,0.53}{\textbf{#1}}}
\newcommand{\ImportTok}[1]{#1}
\newcommand{\InformationTok}[1]{\textcolor[rgb]{0.56,0.35,0.01}{\textbf{\textit{#1}}}}
\newcommand{\KeywordTok}[1]{\textcolor[rgb]{0.13,0.29,0.53}{\textbf{#1}}}
\newcommand{\NormalTok}[1]{#1}
\newcommand{\OperatorTok}[1]{\textcolor[rgb]{0.81,0.36,0.00}{\textbf{#1}}}
\newcommand{\OtherTok}[1]{\textcolor[rgb]{0.56,0.35,0.01}{#1}}
\newcommand{\PreprocessorTok}[1]{\textcolor[rgb]{0.56,0.35,0.01}{\textit{#1}}}
\newcommand{\RegionMarkerTok}[1]{#1}
\newcommand{\SpecialCharTok}[1]{\textcolor[rgb]{0.81,0.36,0.00}{\textbf{#1}}}
\newcommand{\SpecialStringTok}[1]{\textcolor[rgb]{0.31,0.60,0.02}{#1}}
\newcommand{\StringTok}[1]{\textcolor[rgb]{0.31,0.60,0.02}{#1}}
\newcommand{\VariableTok}[1]{\textcolor[rgb]{0.00,0.00,0.00}{#1}}
\newcommand{\VerbatimStringTok}[1]{\textcolor[rgb]{0.31,0.60,0.02}{#1}}
\newcommand{\WarningTok}[1]{\textcolor[rgb]{0.56,0.35,0.01}{\textbf{\textit{#1}}}}
\newcommand{\scauSafeIncludeGraphics}[2][]{%
  \IfFileExists{#2}{\includegraphics[#1]{#2}}{\fbox{\zihao{5} Missing image: #2}}%
}
\newcommand{\scauSubmissionDate}{%%__DATE__}

% Pandoc 对无标题表格生成 \def\LTcaptype{none}，预先定义计数器防止报错
\newcounter{none}

\begin{document}
\pagenumbering{alph}

% 正文基础字体设定为宋体小四，首行缩进两字距
\zihao{-4}\songti
\setlength{\parindent}{2em}

% --------------------封面--------------------
\clearpage
\thispagestyle{empty}
\vfill
\begin{center}
    \scauSafeIncludeGraphics[width=0.8\textwidth]{image/SCAU-LOGO.jpg}
\end{center}
\vfill
\begin{center}
    {\zihao{-0} \songti \textbf{本科毕业论文}}
\end{center}
\vfill
\begin{center}
    {\zihao{2} \heiti \textbf{%%__TITLE_ZH__}}
\end{center}
\vfill
\vfill
\vfill
\begin{center}
    \zihao{-3} \songti
    \renewcommand{\arraystretch}{1.3}
    \begin{tabular}{ll}
        \makebox[4em][s]{学\hfill 院}： & \underline{\makebox[9.5cm]{%%__COLLEGE__}} \\
        \makebox[4em][s]{专\hfill 业}： & \underline{\makebox[9.5cm]{%%__MAJOR__}} \\
        \makebox[4em][s]{姓\hfill 名}： & \underline{\makebox[9.5cm]{%%__AUTHOR_ZH__}} \\
        \makebox[4em][s]{学\hfill 号}： & \underline{\makebox[9.5cm]{%%__STUDENT_ID__}} \\
        \makebox[4em][s]{指导教师}： & \underline{\makebox[3.8cm]{%%__SUPERVISOR__}}\makebox[1.9cm][c]{职称}\underline{\makebox[3.8cm]{%%__SUPERVISOR_TITLE__}} \\
        \makebox[4em][s]{提交日期}： & \underline{\makebox[6.5cm]{%%__DATE__}} \\
    \end{tabular}
\end{center}
\vspace{3em}
\clearpage

% --------------------原创性声明及授权声明--------------------
\clearpage
\thispagestyle{empty}
\vspace{2em}
\begin{center}
    {\zihao{4}\heiti 华南农业大学本科毕业论文（设计）原创性声明}
\end{center}

本人郑重声明：所呈交的毕业论文（设计），是本人在导师的指导下，独立进行研究工作所取得的成果。除文中已经注明引用的内容外，本论文不包含任何其他个人或集体已经发表或撰写过的作品成果。对本文的研究做出重要贡献的个人和集体，均已在文中以明确方式标明。本人完全意识到本声明的法律结果由本人承担。
\vspace{3em}
\begin{center}
    \begin{tabular}{p{0.45\textwidth} p{0.45\textwidth}}
        作者签名：%%__AUTHOR_ZH__ & 日期：%%__DATE__ \\
    \end{tabular}
\end{center}

\vspace{5em}
\begin{center}
    {\zihao{4}\heiti 华南农业大学本科毕业论文（设计）版权使用授权书}
\end{center}

本人完全了解学校有关保留、使用毕业论文（设计）的规定，同意学校保留并向有关部门或机构送交论文（设计）的复印件和电子版，允许论文（设计）被查阅和借阅。
\vspace{3em}
\begin{center}
    \begin{tabular}{p{0.45\textwidth} p{0.45\textwidth}}
        作者签名：%%__AUTHOR_ZH__ & 日期：%%__DATE__ \\
        指导教师签名：%%__SUPERVISOR__ & 日期：%%__DATE__ \\
    \end{tabular}
\end{center}

% --------------------中文摘要-------------------
\clearpage
\pagenumbering{Roman}
\setcounter{page}{1}
\centerline{\zihao{4}\heiti 摘\hspace{4em}要}
\vspace{0.5em}
%%__ABSTRACT_ZH__

\noindent{\zihao{-4}\heiti 关键词：} %%__KEYWORDS_ZH__

% --------------------英文摘要-------------------
\clearpage
\begin{center}
    {\zihao{4}\textbf{%%__TITLE_EN__}}

    {\zihao{-4} %%__AUTHOR_EN__}

    {\zihao{-4} (%%__COLLEGE__, South China Agricultural University, Guangzhou, 510642, China)}
\end{center}
\vspace{-0.89em}
\noindent{\zihao{-4}\textbf{Abstract:}
%%__ABSTRACT_EN__
}

\noindent{\zihao{-4}\textbf{Key words:} %%__KEYWORDS_EN__}

% --------------------符号表-------------------
\clearpage
\centerline{\zihao{4}\heiti 英文缩略词（符号表）}
\begin{center}
    \zihao{-4}
    \begin{tabularx}{\textwidth}{c >{\centering\arraybackslash}X >{\centering\arraybackslash}X}
        \toprule
        英文缩写 & 英文全称 & 中文全称 \\
        \midrule
        LMI & Linear Matrix Inequality & 线性矩阵不等式 \\
        LKF & Lyapunov-Krasovskii Functional & Lyapunov-Krasovskii 泛函 \\
        \bottomrule
    \end{tabularx}
\end{center}

% --------------------目录--------------------
\clearpage
\tableofcontents

% --------------------正文--------------------
\clearpage
\pagenumbering{arabic}
\setcounter{page}{1}

%%__BODY__

% --------------------参考文献--------------------
\clearpage
\phantomsection
\addcontentsline{toc}{section}{参考文献}
\centerline{\zihao{4}\heiti 参考文献}
\vspace{1em}

% 重新定义参考文献环境：取消数字编号，设置悬挂缩进
\makeatletter
\renewcommand{\@biblabel}[1]{} % 取消方括号和数字编号
\renewenvironment{thebibliography}[1]
     {\list{}{\usecounter{enumiv}
              \setlength{\labelwidth}{0pt}
              \setlength{\leftmargin}{2em}    % 整体向右缩进（第二行及之后的缩进量）
              \setlength{\itemindent}{-2em}   % 首行反向缩进，形成悬挂效果
              \setlength{\labelsep}{0pt}      
              \setlength{\itemsep}{0ex}       % 条目之间的垂直间距
              \setlength{\parsep}{0pt}}
      \sloppy
      \clubpenalty4000
      \widowpenalty4000}
     {\endlist}
\makeatother

{\zihao{-4}\songti\rmfamily
\begin{thebibliography}{0}
%%__REFERENCES__
\end{thebibliography}
\par}

% --------------------致谢--------------------
\clearpage
\phantomsection
\addcontentsline{toc}{section}{致谢}
\centerline{\zihao{4}\heiti 致\hspace{4em}谢}
\vspace{1em}
%%__ACKNOWLEDGEMENTS__

% --------------------成绩评定表--------------------
\clearpage
\thispagestyle{empty}
\vspace{1em}

\begin{center}
    {\zihao{3}\heiti 华南农业大学本科毕业论文评分参考标准}
\end{center}

{\zihao{5}\songti
\linespread{1}\selectfont 
\renewcommand{\arraystretch}{1.6} 
\renewcommand{\tabularxcolumn}[1]{m{#1}} 

\noindent 
\begin{tabularx}{\textwidth}{|m{2.2cm}<{\centering}|X|m{1.2cm}<{\centering}|}
    \hline
    \textbf{评价项目} & \multicolumn{1}{c|}{\textbf{评价指标}} & \textbf{分值} \\
    \hline
    \multicolumn{3}{|c|}{\textbf{指导教师评阅}} \\
    \hline
    学习态度\newline 与工作量 & 学习态度认真，工作作风严谨务实；工作量饱满，按期圆满完成规定的任务。 &  \\
    \hline
    资料收集\newline 与分析 & 具有一定的查阅、整理、分析中外文文献的能力；能较好理解课题任务并提出实施方案；查阅文献有一定广泛性，资料收集充分。 &  \\
    \hline
    研究水平\newline 与能力 & 能熟练掌握和运用所学专业基本理论、基本知识和基本技能分析解决相关理论和实际问题；论点鲜明，论据充分；实验设计合理，数据准确可靠，理论分析与计算正确；具有从事科学研究工作或专门技术工作的初步能力。 &  \\
    \hline
    论文撰写\newline 质\qquad 量 & 论文结构严谨，逻辑性强，层次清晰，技术用语准确；行文流畅，语句通顺；论文格式符合规范要求。 &  \\
    \hline
    学术水平\newline 与创新 & 具有一定的学术水平或实践应用价值；对与课题相关的理论或实际问题有较深刻的认识，有新的见解，有一定的创新。 &  \\
    \hline
    
    \multicolumn{3}{|c|}{\textbf{评阅教师评阅}} \\
    \hline
    选 题 质 量 & 选题符合专业培养目标要求；与科学研究、生产实际或社会实践问题紧密结合；工作量饱满，难度适中，综合训练强。 &  \\
    \hline
    资料收集\newline 与分析 & 具有一定的查阅、整理、分析中外文文献的能力；能较好理解课题任务并提出实施方案；查阅文献有一定广泛性，资料收集充分。 &  \\
    \hline
    研究水平\newline 与能力 & 能熟练掌握和运用所学专业基本理论、基本知识和基本技能分析解决相关理论和实际问题；论点鲜明，论据充分；实验设计合理，数据准确可靠，理论分析与计算正确；具有从事科学研究工作或专门技术工作的初步能力。 &  \\
    \hline
    论文撰写\newline 质\qquad 量 & 论文结构严谨，逻辑性强，层次清晰，技术用语准确；行文流畅，语句通顺；论文格式符合规范要求。 &  \\
    \hline
    学术水平\newline 与创新 & 具有一定的学术水平或实践应用价值；对与课题相关的理论或实际问题有较深刻的认识，有新的见解，有一定的创新。 &  \\
    \hline
    
    \multicolumn{3}{|c|}{\textbf{答辩小组评阅}} \\
    \hline
    选 题 质 量 & 选题符合专业培养目标要求；与科学研究、生产实际或社会实践问题紧密结合；工作量饱满，难度适中，综合训练强。 &  \\
    \hline
    研究水平\newline 与能力 & 能熟练掌握和运用所学专业基本理论、基本知识和基本技能分析解决相关理论和实际问题；论点鲜明，论据充分；实验设计合理，数据准确可靠，理论分析与计算正确；具有从事科学研究工作或专门技术工作的初步能力。 &  \\
    \hline
    论文撰写\newline 质\qquad 量 & 论文结构严谨，逻辑性强，层次清晰，技术用语准确；行文流畅，语句通顺；论文格式符合规范要求。 &  \\
    \hline
    学术水平\newline 与创新 & 具有一定的学术水平或实践应用价值；对与课题相关的理论或实际问题有较深刻的认识，有新的见解，有一定的创新。 &  \\
    \hline
    答\qquad 辩 & 能简明扼要阐述论文主要内容，思路清晰，语言表达准确、顺畅；回答问题有理论根据，基本概念清楚，逻辑性强，对主要问题回答准确、有深度。 &  \\
    \hline
\end{tabularx}
\par}


\end{document}
